% Options for packages loaded elsewhere
\PassOptionsToPackage{unicode}{hyperref}
\PassOptionsToPackage{hyphens}{url}
%
\documentclass[
]{article}
\usepackage{amsmath,amssymb}
\usepackage{iftex}
\ifPDFTeX
  \usepackage[T1]{fontenc}
  \usepackage[utf8]{inputenc}
  \usepackage{textcomp} % provide euro and other symbols
\else % if luatex or xetex
  \usepackage{unicode-math} % this also loads fontspec
  \defaultfontfeatures{Scale=MatchLowercase}
  \defaultfontfeatures[\rmfamily]{Ligatures=TeX,Scale=1}
\fi
\usepackage{lmodern}
\ifPDFTeX\else
  % xetex/luatex font selection
\fi
% Use upquote if available, for straight quotes in verbatim environments
\IfFileExists{upquote.sty}{\usepackage{upquote}}{}
\IfFileExists{microtype.sty}{% use microtype if available
  \usepackage[]{microtype}
  \UseMicrotypeSet[protrusion]{basicmath} % disable protrusion for tt fonts
}{}
\makeatletter
\@ifundefined{KOMAClassName}{% if non-KOMA class
  \IfFileExists{parskip.sty}{%
    \usepackage{parskip}
  }{% else
    \setlength{\parindent}{0pt}
    \setlength{\parskip}{6pt plus 2pt minus 1pt}}
}{% if KOMA class
  \KOMAoptions{parskip=half}}
\makeatother
\usepackage{xcolor}
\usepackage{color}
\usepackage{fancyvrb}
\newcommand{\VerbBar}{|}
\newcommand{\VERB}{\Verb[commandchars=\\\{\}]}
\DefineVerbatimEnvironment{Highlighting}{Verbatim}{commandchars=\\\{\}}
% Add ',fontsize=\small' for more characters per line
\newenvironment{Shaded}{}{}
\newcommand{\AlertTok}[1]{\textcolor[rgb]{1.00,0.00,0.00}{\textbf{#1}}}
\newcommand{\AnnotationTok}[1]{\textcolor[rgb]{0.38,0.63,0.69}{\textbf{\textit{#1}}}}
\newcommand{\AttributeTok}[1]{\textcolor[rgb]{0.49,0.56,0.16}{#1}}
\newcommand{\BaseNTok}[1]{\textcolor[rgb]{0.25,0.63,0.44}{#1}}
\newcommand{\BuiltInTok}[1]{\textcolor[rgb]{0.00,0.50,0.00}{#1}}
\newcommand{\CharTok}[1]{\textcolor[rgb]{0.25,0.44,0.63}{#1}}
\newcommand{\CommentTok}[1]{\textcolor[rgb]{0.38,0.63,0.69}{\textit{#1}}}
\newcommand{\CommentVarTok}[1]{\textcolor[rgb]{0.38,0.63,0.69}{\textbf{\textit{#1}}}}
\newcommand{\ConstantTok}[1]{\textcolor[rgb]{0.53,0.00,0.00}{#1}}
\newcommand{\ControlFlowTok}[1]{\textcolor[rgb]{0.00,0.44,0.13}{\textbf{#1}}}
\newcommand{\DataTypeTok}[1]{\textcolor[rgb]{0.56,0.13,0.00}{#1}}
\newcommand{\DecValTok}[1]{\textcolor[rgb]{0.25,0.63,0.44}{#1}}
\newcommand{\DocumentationTok}[1]{\textcolor[rgb]{0.73,0.13,0.13}{\textit{#1}}}
\newcommand{\ErrorTok}[1]{\textcolor[rgb]{1.00,0.00,0.00}{\textbf{#1}}}
\newcommand{\ExtensionTok}[1]{#1}
\newcommand{\FloatTok}[1]{\textcolor[rgb]{0.25,0.63,0.44}{#1}}
\newcommand{\FunctionTok}[1]{\textcolor[rgb]{0.02,0.16,0.49}{#1}}
\newcommand{\ImportTok}[1]{\textcolor[rgb]{0.00,0.50,0.00}{\textbf{#1}}}
\newcommand{\InformationTok}[1]{\textcolor[rgb]{0.38,0.63,0.69}{\textbf{\textit{#1}}}}
\newcommand{\KeywordTok}[1]{\textcolor[rgb]{0.00,0.44,0.13}{\textbf{#1}}}
\newcommand{\NormalTok}[1]{#1}
\newcommand{\OperatorTok}[1]{\textcolor[rgb]{0.40,0.40,0.40}{#1}}
\newcommand{\OtherTok}[1]{\textcolor[rgb]{0.00,0.44,0.13}{#1}}
\newcommand{\PreprocessorTok}[1]{\textcolor[rgb]{0.74,0.48,0.00}{#1}}
\newcommand{\RegionMarkerTok}[1]{#1}
\newcommand{\SpecialCharTok}[1]{\textcolor[rgb]{0.25,0.44,0.63}{#1}}
\newcommand{\SpecialStringTok}[1]{\textcolor[rgb]{0.73,0.40,0.53}{#1}}
\newcommand{\StringTok}[1]{\textcolor[rgb]{0.25,0.44,0.63}{#1}}
\newcommand{\VariableTok}[1]{\textcolor[rgb]{0.10,0.09,0.49}{#1}}
\newcommand{\VerbatimStringTok}[1]{\textcolor[rgb]{0.25,0.44,0.63}{#1}}
\newcommand{\WarningTok}[1]{\textcolor[rgb]{0.38,0.63,0.69}{\textbf{\textit{#1}}}}
\setlength{\emergencystretch}{3em} % prevent overfull lines
\providecommand{\tightlist}{%
  \setlength{\itemsep}{0pt}\setlength{\parskip}{0pt}}
\setcounter{secnumdepth}{-\maxdimen} % remove section numbering
\ifLuaTeX
  \usepackage{selnolig}  % disable illegal ligatures
\fi
\usepackage{bookmark}
\IfFileExists{xurl.sty}{\usepackage{xurl}}{} % add URL line breaks if available
\urlstyle{same}
\hypersetup{
  hidelinks,
  pdfcreator={LaTeX via pandoc}}

\author{}
\date{}

\begin{document}

\section{Legacy Gear-Spin Direction
Reference}\label{legacy-gear-spin-direction-reference}

\subsection{Scope and Disposition}\label{scope-and-disposition}

Issue
\href{https://github.com/D-sorganization/UpstreamDrift/issues/12006}{\#12006}
corrects the horizontal spin sign in the public legacy impact helper.
The matched-delivery study uses coupled contact and does not call this
helper. Historical study outputs and their source hashes remain
unchanged. This is a separate research reference; its editable
\href{LEGACY_GEAR_SIGN_REFERENCE.tex}{LaTeX Source} is not part of the
canonical engineering manual.

\subsection{Frames, Inputs, and
Equations}\label{frames-inputs-and-equations}

For the documented right-handed convention, X points forward, Z upward,
and Y left. The two contact coordinates are toe-positive horizontal
offset h and high-positive vertical offset v, in metres from face
center. V is clubhead speed in m/s, g is a dimensionless gear factor,
and \(k_h/k_v\) are heuristic scales in \(\mathrm{rad\,m^{-2}}\). With
unit upward vector \(\mathbf u\) and horizontal face-normal
perpendicular
\(\mathbf t = (\mathbf u \times \mathbf n)/|\mathbf u \times \mathbf n|\),
the additive angular velocity is

\[\Delta\boldsymbol\omega = g h V k_h\,\mathbf u + g v V k_v\,\mathbf t.\]

For a forward face normal, t is +Y. The native Magnus direction is
angular velocity cross ball velocity. Thus positive h adds +Z spin and
produces a leftward contribution for a +X shot. Negative h reverses it.
Positive v adds topspin about +Y and reduces negative-Y backspin. The
vertical term, scale, contact model and native flight equations are
unchanged by this repair.

The old horizontal term had the opposite sign. At h = 0.00635 m, V = 45
m/s, g = 0.5 and \(k_h = 350\,\mathrm{rad\,m^{-2}}\), the corrected
contribution is +50.00625 rad/s about Z; the old code returned
\(-50.00625\,\mathrm{rad/s}\). This is a frame/sign correction, not a
recalibration or a change in the magnitude model.

\subsection{Acceptance Evidence and
Reproduction}\label{acceptance-evidence-and-reproduction}

Nine new tests failed before the production change: toe/heel Magnus
direction at three headings, public solver spin increments, and
public-solver-to-native landing direction. The native pre-fix test
placed the heel shot at +11.293335 m laterally while the centered shot
stayed at zero, demonstrating the reversed curvature. After the sign
correction, those tests and the existing impact regressions pass. An
additional boundary regression replaces the legacy helper with a
function that raises and verifies the matched contact pathway still
produces a finite off-center landing.

The old assertions that called negative-Z toe spin a draw were corrected
to agree with the independently tested force direction. Their magnitude
intervals were reflected, not widened. Astra independently reviewed the
sign and ran the original nine new tests. The final relevant suite
passes 58 tests, including both native flight tests. The
\href{LEGACY_GEAR_SIGN_VALIDATION.json}{Validation Receipt} binds source
hashes and the \href{LEGACY_GEAR_SIGN_TESTS.txt}{Final Test Log}; the
\href{LEGACY_GEAR_SIGN_RED.txt}{Pre-Fix Log} preserves the nine failing
regressions. LaTeX compilation remains unverified because the built-in
compiler could not retrieve its uncached bundle; editable source is
preserved.

\begin{Shaded}
\begin{Highlighting}[]
\VariableTok{OPENBLAS\_NUM\_THREADS}\OperatorTok{=}\NormalTok{1 }\ExtensionTok{/tmp/shot{-}pattern{-}validation{-}venv{-}20261009/bin/python} \AttributeTok{{-}m}\NormalTok{ pytest tests/unit/physics/test\_legacy\_gear\_flight\_direction.py tests/unit/physics/test\_impact\_physics\_value\_assertions.py tests/unit/shared\_python/test\_impact\_model.py tests/unit/physics/test\_impact\_friction\_axis\_and\_gear\_offset.py }\AttributeTok{{-}o}\NormalTok{ addopts=}\StringTok{\textquotesingle{}\textquotesingle{}} \AttributeTok{{-}q}
\end{Highlighting}
\end{Shaded}

Native flight tests skip explicitly if the required kernel cannot be
imported; the qualified local run must execute them. Pure vector and
public solver tests do not depend on a native flight installation.

\subsection{Rigid-Contact Gear Law for the Shot
Card}\label{rigid-contact-gear-law-for-the-shot-card}

Issue \#15443 adds an opt-in replacement for the heuristic. When
\texttt{ImpactParameters.rigid\_gear\_head} is set,
\texttt{ImpactSolverAPI} takes the gear increment from
\texttt{physics/impact\_model/rigid\_gear.py} instead:

\[\Delta\boldsymbol\omega = \boldsymbol\omega_b(b\,\mathbf n - h\,\mathbf e_h + v\,\mathbf e_u) - \boldsymbol\omega_b(b\,\mathbf n),\]

where \(\boldsymbol\omega_b(\mathbf r)\) is the ball spin from the shared
finite-inertia contact solver (\texttt{physics/rigid\_contact.py}). The
solver uses the contact lever \(\mathbf r\) from the head CG, the head
inertia tensor
\(\mathbf R\,\mathrm{diag}(I_n, I_h, I_u)\,\mathbf R^\top\) in the face
frame (normal \(\mathbf n\), heel
\(\mathbf e_h = \mathbf u \times \mathbf n\), face-up \(\mathbf e_u\)),
CG depth \(b\), head mass, COR and friction. The face is flat. Only the
spin increment is applied; the offset's change to ball speed and launch
direction is not.

Head values per club are settings
(\texttt{impact\_parameters/club\_impact\_settings.json},
\texttt{rigid\_gear\_heads}). They are the phase-two matched-delivery
research assumptions, not measured heads. Only the shot card and export
flight path (\texttt{flight\_report.flight\_for\_impact}) opt in. Every
other caller keeps the legacy term and its historical outputs.

Evidence comes from
\texttt{tests/unit/physics/test\_rigid\_gear\_effect\_15443.py}:

\begin{itemize}
\item A centred strike gives zero increment.
\item Toe and heel strikes are antisymmetric, with the legacy sign.
\item A CG in the face plane gives no side spin.
\item A higher yaw inertia gives less gear.
\end{itemize}

For a measured driver heel strike (h = \(-3.4\) mm, v = \(+1.4\) mm,
51.4 m/s), the law gives \(-143\) rpm of side spin against the
heuristic's \(-291\) rpm. On the matched driver swing, this moves the
spin axis from 19.8° to 15.3° (D-plane alone: 10.3°). On the matched
7-iron (CG depth 10 mm), the gear tilt falls from 1.5° to 0.15°.

This is a software regression on a reviewed contact law, not an
equipment calibration: no launch-monitor comparison qualifies these head
values. Face bulge and roll are the next slice.

\subsection{Limitations}\label{limitations}

The legacy model adds spin without solving the coupled club/ball
rotational collision. Its coefficients have no player or equipment
calibration in this study, and its angular momentum and energy
consistency are not established by this sign repair. It assumes an
approximately forward face normal and a global vertical spin axis; it is
not an arbitrary-orientation contact law. No accuracy claim,
left-handed-delivery qualification, or universal net-curvature claim
follows from these tests. Bulge, roll and other launch/spin
contributions can change net curvature in a fuller model. Use the
independently qualified coupled contact pathway for the matched-delivery
research.

\end{document}
